\documentclass[12pt, a4paper]{article}

\usepackage[utf8]{inputenc}
\usepackage[T1]{fontenc}
\usepackage[french]{babel}
\usepackage{amsmath, amssymb, amsthm}
\usepackage{geometry}
\usepackage{graphicx}
\usepackage{hyperref}
\usepackage{booktabs}
\usepackage{array}
\usepackage{tikz}
\usetikzlibrary{arrows.meta, positioning, calc}
\usepackage{xcolor}
\usepackage{listings}
\usepackage{float}

\geometry{margin=2.5cm}

\definecolor{bleu}{RGB}{30, 100, 200}
\definecolor{rouge}{RGB}{200, 50, 50}
\definecolor{vert}{RGB}{30, 150, 80}
\definecolor{gris}{RGB}{100, 100, 100}

\hypersetup{
    colorlinks=true,
    linkcolor=bleu,
    urlcolor=bleu,
    citecolor=bleu
}

\newtheorem{definition}{Définition}
\newtheorem{propriete}{Propriété}
\newtheorem{theoreme}{Théorème}
\newcommand{\dist}[2]{d\!\left(#1,\,#2\right)}

\title{
    \textbf{Arbre de Steiner euclidien pour 4 points}\\[0.4em]
    \large Fonctionnement complet de l'algorithme —\\
    du clic utilisateur au résultat affiché
}
\author{Projet Algorithmes 2}
\date{\today}

% ------------------------------------------------------------------
\begin{document}
\maketitle
\tableofcontents
\newpage

% ==================================================================
\section{Vue d'ensemble du système}
% ==================================================================

L'application est composée de deux parties qui communiquent via HTTP :

\begin{center}
\begin{tikzpicture}[
    box/.style={draw, rounded corners, minimum width=3.2cm, minimum height=1cm,
                fill=#1!10, draw=#1!60, font=\small},
    arr/.style={-{Stealth}, thick}
]
    \node[box=bleu] (front) {Frontend Angular};
    \node[box=rouge, right=3.5cm of front] (back) {Backend Spring Boot};
    \node[box=vert, right=3.5cm of back] (algo) {SteinerTreeService};

    \draw[arr] (front) -- node[above, font=\footnotesize]{POST /api/steiner/solve} (back);
    \draw[arr] (back)  -- node[above, font=\footnotesize]{solve(points)} (algo);
    \draw[arr] (algo)  to[bend right=30] node[below, font=\footnotesize]{SteinerResult} (back);
    \draw[arr] (back)  to[bend right=30] node[below, font=\footnotesize]{JSON} (front);
\end{tikzpicture}
\end{center}

\bigskip

Lorsque l'utilisateur place exactement \textbf{4 points} sur le canvas et clique sur
\emph{Calculer l'arbre de Steiner}, la séquence suivante se déclenche :

\begin{enumerate}
    \item Le composant Angular \texttt{AppComponent} collecte les coordonnées $(x_i, y_i)$.
    \item \texttt{SteinerService.solve()} envoie un \texttt{POST} au backend avec la liste JSON des 4 points.
    \item \texttt{SteinerController} reçoit la requête et délègue à \texttt{SteinerTreeService.solve()}.
    \item Le service détecte $n = 4$ et appelle \texttt{solveForFourPoints()}.
    \item Le résultat (arêtes + points de Steiner) est renvoyé en JSON et affiché sur le canvas.
\end{enumerate}

% ==================================================================
\section{Problème résolu — définition mathématique}
% ==================================================================

\begin{definition}[Arbre de Steiner euclidien]
Étant donnés $n$ points terminaux $T_1, \ldots, T_n \in \mathbb{R}^2$, l'arbre de Steiner
euclidien est l'arbre connexe de longueur minimale qui les relie tous, en autorisant
l'ajout de nouveaux points intermédiaires appelés \textbf{points de Steiner}.
\end{definition}

\begin{propriete}[Angles à 120°]
Dans tout arbre de Steiner euclidien optimal, les arêtes se rejoignent à exactement
$120°$ en chaque point de Steiner.
\end{propriete}

\begin{proof}
Soit $S$ un point de Steiner de degré 3 relié à $A$, $B$, $C$. On minimise
$f(S) = \dist{S}{A} + \dist{S}{B} + \dist{S}{C}$.
Le gradient est $\nabla f(S) = \dfrac{S-A}{\|S-A\|} + \dfrac{S-B}{\|S-B\|} + \dfrac{S-C}{\|S-C\|}$.
À l'optimum, $\nabla f = \mathbf{0}$, donc les trois vecteurs unitaires sortants
se compensent, ce qui impose des angles de $120°$ entre chaque paire.
\end{proof}

\begin{propriete}[Nombre maximal de points de Steiner]
Pour $n$ terminaux, un arbre de Steiner optimal contient \textbf{au plus $n-2$}
points de Steiner.
\end{propriete}

\begin{proof}
Un arbre sur $n + k$ nœuds possède exactement $n + k - 1$ arêtes.
Chaque point de Steiner est de degré~3 (sinon on peut le supprimer ou le scinder
pour raccourcir l'arbre).
La somme des degrés vaut $2(n + k - 1)$.
Les $n$ terminaux ont degré $\geq 1$, les $k$ points de Steiner contribuent $3k$ :
\[
    2(n + k - 1) \;\geq\; n + 3k
    \;\Longrightarrow\; n \;\geq\; k + 2
    \;\Longrightarrow\; k \;\leq\; n - 2.
\]
Pour $n = 4$ : \textbf{au plus 2 points de Steiner}.
\end{proof}

% ==================================================================
\section{Stratégie globale pour 4 points}
% ==================================================================

Puisqu'il y a au plus 2 points de Steiner, on distingue trois familles de topologies :

\begin{center}
\begin{tabular}{clcc}
\toprule
\textbf{Famille} & \textbf{Structure} & \textbf{Nb. topologies} & \textbf{Méthode} \\
\midrule
0 Steiner & Arbre couvrant minimal (MST) & 1 & Algorithme de Prim \\
1 Steiner & Triplet + terminal isolé & 12 & Point de Fermat-Torricelli \\
2 Steiner & Partition en deux paires & 3 & Weiszfeld alterné \\
\midrule
& \textbf{Total} & \textbf{16} & \\
\bottomrule
\end{tabular}
\end{center}

\medskip

L'algorithme évalue toutes les topologies valides et retourne celle de longueur minimale.

% ==================================================================
\section{Famille 0 — Arbre couvrant minimal (MST)}
% ==================================================================

\subsection{Définition}

L'arbre couvrant minimal connecte les 4 terminaux sans point de Steiner supplémentaire.
Il constitue une \textbf{borne supérieure} : toute solution avec des points de Steiner
doit être strictement plus courte pour être retenue.

\subsection{Algorithme de Prim}

On initialise l'arbre avec $T_1$ et on ajoute itérativement l'arête de poids minimum
vers un sommet non encore visité.

\begin{align}
&\text{inMST} = \{T_1\},\quad \text{parent}[i] = -1,\quad \text{minDist}[i] = +\infty \\
&\text{minDist}[1] = 0 \nonumber\\[4pt]
&\textbf{Répéter 3 fois :} \nonumber\\
&\quad u \leftarrow \arg\min_{i \notin \text{inMST}} \text{minDist}[i] \nonumber\\
&\quad \text{inMST} \leftarrow \text{inMST} \cup \{u\} \nonumber\\
&\quad \forall v \notin \text{inMST} :
       \quad \text{si } \dist{T_u}{T_v} < \text{minDist}[v]
       \;\text{ alors }\; \text{minDist}[v] \leftarrow \dist{T_u}{T_v},\;
       \text{parent}[v] \leftarrow u \nonumber
\end{align}

\textbf{Complexité :} $O(n^2)$, négligeable pour $n = 4$.

% ==================================================================
\section{Famille 1 — Un point de Steiner (12 topologies)}
% ==================================================================

\subsection{Structure}

On choisit un sous-ensemble de 3 terminaux (\emph{triplet}) et un terminal restant
(\emph{lone}). Le point de Steiner $S$ est le point de Fermat-Torricelli du triplet.
Le terminal \emph{lone} se raccorde directement à l'un des 3 terminaux du triplet
(3 choix).

\medskip

Il y a $\binom{4}{3} = 4$ triplets possibles, chacun donnant 3 attaches, soit
$4 \times 3 = 12$ topologies.

\begin{center}
\begin{tikzpicture}[
    pt/.style={circle, fill=bleu!80, minimum size=6pt, inner sep=0pt},
    st/.style={circle, fill=rouge!80, minimum size=6pt, inner sep=0pt},
    edge/.style={thick, gris}
]
    \node[pt, label=above left:$T_1$] (T1) at (0, 1.5) {};
    \node[pt, label=above right:$T_2$] (T2) at (2, 1.5) {};
    \node[pt, label=below:$T_3$] (T3) at (1, 0) {};
    \node[pt, label=right:$T_4$ (lone)] (T4) at (3.5, 0.7) {};
    \node[st, label=above:$S$] (S) at (1, 0.9) {};

    \draw[edge] (S) -- (T1);
    \draw[edge] (S) -- (T2);
    \draw[edge] (S) -- (T3);
    \draw[edge, dashed] (T4) -- (T2) node[midway, above, font=\footnotesize]{attache};
\end{tikzpicture}
\end{center}

\subsection{Point de Fermat-Torricelli}

\begin{definition}[Point de Fermat-Torricelli]
Étant donnés trois points $A$, $B$, $C$, le point de Fermat-Torricelli est le point
$F$ qui minimise :
\[
    f(F) = \dist{F}{A} + \dist{F}{B} + \dist{F}{C}.
\]
\end{definition}

\begin{propriete}[Cas du sommet obtus]
Si l'angle en un sommet du triangle $ABC$ est supérieur ou égal à $120°$, alors
$F$ coïncide avec ce sommet.
\end{propriete}

\begin{proof}
Notons $\alpha$ l'angle en $A$. On a $f(A) = \dist{A}{B} + \dist{A}{C}$.
Déplacer $F$ au voisinage de $A$ ne peut qu'augmenter $f$, car le gradient de
$\dist{F}{B} + \dist{F}{C}$ en $F = A$ vaut $\dfrac{A-B}{\|A-B\|} + \dfrac{A-C}{\|A-C\|}$,
dont la norme est $2\cos(\alpha/2) \leq 1$ pour $\alpha \geq 120°$.
\end{proof}

\subsection{Algorithme de Weiszfeld (1937)}

Lorsqu'aucun angle n'atteint $120°$, le point de Fermat est calculé par itérations :

\begin{equation}
\label{eq:weiszfeld}
    F^{(k+1)} \;=\;
    \frac{\displaystyle\frac{A}{\dist{F^{(k)}}{A}}
        + \frac{B}{\dist{F^{(k)}}{B}}
        + \frac{C}{\dist{F^{(k)}}{C}}}
         {\displaystyle\frac{1}{\dist{F^{(k)}}{A}}
        + \frac{1}{\dist{F^{(k)}}{B}}
        + \frac{1}{\dist{F^{(k)}}{C}}}
\end{equation}

\textbf{Interprétation :} $F^{(k+1)}$ est la \emph{moyenne pondérée} de $A$, $B$, $C$
avec des poids $w_i = 1/\dist{F^{(k)}}{P_i}$ (les points éloignés ont moins de poids).

\medskip

\textbf{Convergence :} La suite $(F^{(k)})$ converge vers le minimum global si $F^{(0)}$
n'est pas confondu avec l'un des $P_i$. On démarre depuis le centroïde
$F^{(0)} = (A + B + C)/3$.

\medskip

\textbf{Critère d'arrêt :}
\[
    \|F^{(k+1)} - F^{(k)}\| \;<\; \varepsilon = 10^{-10}
    \quad \text{(ou au bout de 50\,000 itérations)}
\]

\subsection{Cas particuliers dans l'implémentation}

\begin{enumerate}
    \item \textbf{Triangle dégénéré (colinéaire) :}
          $|\overrightarrow{AB} \times \overrightarrow{AC}| < 10^{-10}$ $\Rightarrow$
          retourne \texttt{null}, topologie ignorée.
    \item \textbf{Angle $\geq 120°$ en $A$ :} $\cos(\angle BAC) \leq -\tfrac{1}{2}$
          $\Rightarrow$ $F = A$, pas de point de Steiner utile.
    \item \textbf{$F$ proche d'un sommet :} $\dist{F}{T_i} < 10^{-8}$
          $\Rightarrow$ pas de point de Steiner ajouté au résultat.
\end{enumerate}

% ==================================================================
\section{Famille 2 — Deux points de Steiner (3 topologies)}
% ==================================================================

\subsection{Structure}

On partitionne les 4 terminaux en deux paires :
\[
    \{T_a, T_b\} \;|\; \{T_c, T_d\}.
\]
Les trois partitions non ordonnées sont :
$\{0,1\}|\{2,3\}$,\quad $\{0,2\}|\{1,3\}$,\quad $\{0,3\}|\{1,2\}$.

La topologie est :
\[
    S_1 \;\leftrightarrow\; T_a,\; T_b,\; S_2
    \qquad \text{et} \qquad
    S_2 \;\leftrightarrow\; T_c,\; T_d,\; S_1.
\]

\begin{center}
\begin{tikzpicture}[
    pt/.style={circle, fill=bleu!80, minimum size=6pt, inner sep=0pt},
    st/.style={circle, fill=rouge!80, minimum size=6pt, inner sep=0pt},
    edge/.style={thick, gris}
]
    \node[pt, label=above left:$T_a$] (Ta) at (0, 1.5) {};
    \node[pt, label=below left:$T_b$] (Tb) at (0, 0) {};
    \node[pt, label=above right:$T_c$] (Tc) at (4, 1.5) {};
    \node[pt, label=below right:$T_d$] (Td) at (4, 0) {};
    \node[st, label=above:$S_1$] (S1) at (1.3, 0.75) {};
    \node[st, label=above:$S_2$] (S2) at (2.7, 0.75) {};

    \draw[edge] (S1) -- (Ta);
    \draw[edge] (S1) -- (Tb);
    \draw[edge] (S1) -- (S2);
    \draw[edge] (S2) -- (Tc);
    \draw[edge] (S2) -- (Td);
\end{tikzpicture}
\end{center}

\subsection{Problème d'optimisation}

On minimise la longueur totale de l'arbre en fonction de $(S_1, S_2) \in \mathbb{R}^4$ :
\begin{equation}
\label{eq:objective}
    \min_{S_1, S_2} f(S_1, S_2) \;=\;
    \dist{S_1}{T_a} + \dist{S_1}{T_b} + \dist{S_1}{S_2}
    + \dist{S_2}{T_c} + \dist{S_2}{T_d}.
\end{equation}

Ce problème est \textbf{convexe} (somme de normes), donc tout minimum local est global.

\subsection{Weiszfeld alterné (coordinate descent)}

On alterne la mise à jour de $S_1$ et $S_2$, chacune étant un sous-problème de
Weber (point de Fermat de 3 points) :

\bigskip

\textbf{Étape 1 — Mise à jour de $S_1$ en fixant $S_2$}

$S_1$ est le point de Fermat de $\{T_a, T_b, S_2^{(k)}\}$ :
\begin{equation}
\label{eq:update_s1}
    S_1^{(k+1)} \;=\;
    \frac{\dfrac{T_a}{\dist{S_1^{(k)}}{T_a}}
        + \dfrac{T_b}{\dist{S_1^{(k)}}{T_b}}
        + \dfrac{S_2^{(k)}}{\dist{S_1^{(k)}}{S_2^{(k)}}}}
         {\dfrac{1}{\dist{S_1^{(k)}}{T_a}}
        + \dfrac{1}{\dist{S_1^{(k)}}{T_b}}
        + \dfrac{1}{\dist{S_1^{(k)}}{S_2^{(k)}}}}
\end{equation}

\textbf{Étape 2 — Mise à jour de $S_2$ en fixant $S_1$}

\begin{equation}
\label{eq:update_s2}
    S_2^{(k+1)} \;=\;
    \frac{\dfrac{T_c}{\dist{S_2^{(k)}}{T_c}}
        + \dfrac{T_d}{\dist{S_2^{(k)}}{T_d}}
        + \dfrac{S_1^{(k+1)}}{\dist{S_2^{(k)}}{S_1^{(k+1)}}}}
         {\dfrac{1}{\dist{S_2^{(k)}}{T_c}}
        + \dfrac{1}{\dist{S_2^{(k)}}{T_d}}
        + \dfrac{1}{\dist{S_2^{(k)}}{S_1^{(k+1)}}}}
\end{equation}

\textbf{Critère d'arrêt :}
\[
    \max\!\left(
        \|S_1^{(k+1)} - S_1^{(k)}\|,\;
        \|S_2^{(k+1)} - S_2^{(k)}\|
    \right) \;<\; 10^{-10}
    \quad \text{(ou 50\,000 itérations)}
\]

\subsection{Initialisations multiples}

Pour éviter les mauvais points de départ, on lance l'optimisation depuis 4 positions
initiales différentes et on conserve la meilleure :

\begin{center}
\begin{tabular}{ll}
\toprule
\textbf{Init} & \textbf{Position initiale $(S_1^{(0)}, S_2^{(0)})$} \\
\midrule
1 & Milieux $\overline{T_a T_b}$ et $\overline{T_c T_d}$ \\
2 & Milieux décalés de $\tfrac{1}{3}$ vers le centroïde global \\
3 & Barycentre de $\{T_a, T_b, \mathrm{centroïde}\}$ et de $\{T_c, T_d, \mathrm{centroïde}\}$ \\
4 & $0{,}65\,T_a + 0{,}35\,\mathrm{centroïde}$ et $0{,}65\,T_c + 0{,}35\,\mathrm{centroïde}$ \\
\bottomrule
\end{tabular}
\end{center}

\subsection{Validation géométrique}

Après optimisation, on rejette la solution si :

\begin{enumerate}
    \item $\dist{S_1}{S_2} < \varepsilon$ \quad (les deux Steiner sont confondus),
    \item $\dist{S_i}{T_j} < \varepsilon$ pour un terminal quelconque \quad (Steiner sur terminal),
    \item Les angles en $S_1$ ou $S_2$ s'écartent de $120°$ de plus de $10°$.
\end{enumerate}

\noindent La valeur $\varepsilon = \max(10^{-8},\; 10^{-4} \times \Delta)$ où
$\Delta = \max_{i < j} \dist{T_i}{T_j}$ est l'échelle du problème.

\textbf{Vérification des angles :} pour le point $S_1$ relié à $T_a$, $T_b$, $S_2$,
on calcule les trois angles :
\[
    \theta_{ab} = \angle(T_a, S_1, T_b),\quad
    \theta_{aS} = \angle(T_a, S_1, S_2),\quad
    \theta_{bS} = \angle(T_b, S_1, S_2),
\]
où $\angle(P, Q, R) = \arccos\!\left(\dfrac{\overrightarrow{QP} \cdot \overrightarrow{QR}}
{\|\overrightarrow{QP}\|\,\|\overrightarrow{QR}\|}\right)$.

La topologie est valide si et seulement si chacun de ces angles est dans $[110°, 130°]$.

% ==================================================================
\section{Sélection de la meilleure topologie}
% ==================================================================

Après avoir évalué toutes les topologies valides, l'algorithme sélectionne celle
de longueur totale minimale :

\[
    \text{best} \;=\; \arg\min_{\text{topologie valide}} \ell(\text{topologie}),
\]

où $\ell$ est la somme des longueurs de toutes les arêtes.

\medskip

\textbf{Note :} la topologie MST (0 Steiner) est toujours valide et sert de borne
supérieure. L'algorithme retourne donc toujours un résultat.

% ==================================================================
\section{Récapitulatif du flux de données}
% ==================================================================

\begin{center}
\begin{tikzpicture}[
    step/.style={draw, rounded corners, fill=#1!10, draw=#1!60,
                 minimum width=6cm, minimum height=0.8cm, font=\small},
    arr/.style={-{Stealth}, thick}
]
    \node[step=bleu] (A) {Utilisateur clique sur le canvas (4 points)};
    \node[step=bleu, below=0.4cm of A] (B) {AppComponent.onPointAdded()};
    \node[step=bleu, below=0.4cm of B] (C) {SteinerService.solve() — POST JSON};
    \node[step=rouge, below=0.4cm of C] (D) {SteinerController.solve()};
    \node[step=rouge, below=0.4cm of D] (E) {SteinerTreeService.solveForFourPoints()};
    \node[step=vert, below=0.4cm of E] (F1) {Évaluation MST (Prim)};
    \node[step=vert, below=0.4cm of F1] (F2) {12 topologies 1-Steiner (Fermat-Weiszfeld)};
    \node[step=vert, below=0.4cm of F2] (F3) {3 topologies 2-Steiner (Weiszfeld alterné)};
    \node[step=rouge, below=0.4cm of F3] (G) {Sélection du minimum — SteinerResult};
    \node[step=bleu, below=0.4cm of G] (H) {Affichage sur le canvas Angular};

    \foreach \x/\y in {A/B, B/C, C/D, D/E, E/F1, F1/F2, F2/F3, F3/G, G/H}
        \draw[arr] (\x) -- (\y);
\end{tikzpicture}
\end{center}

% ==================================================================
\section{Exemple numérique — carré unité}
% ==================================================================

Soit les 4 terminaux $T_1=(0,0)$, $T_2=(1,0)$, $T_3=(1,1)$, $T_4=(0,1)$.

\subsection{MST}
La longueur du MST vaut $3$ (trois côtés du carré).

\subsection{1 Steiner — triplet $\{T_1, T_2, T_4\}$, lone $= T_3$}
Les angles intérieurs du triangle rectangle isocèle valent $90°, 45°, 45°$.
Aucun n'atteint $120°$, donc on calcule le point de Fermat par Weiszfeld.
La convergence donne $S \approx (0{,}5, 0{,}211)$.
Longueur totale $\approx 2{,}866$.

\subsection{2 Steiner — partition $\{T_1, T_4\}|\{T_2, T_3\}$}
Par symétrie, $S_1 = (0, 0{,}5)$ et $S_2 = (1, 0{,}5)$ ne conviennent pas
(angles de $90°$ seulement). L'optimisation converge vers
$S_1 \approx (0{,}211, 0{,}5)$ et $S_2 \approx (0{,}789, 0{,}5)$.
Longueur totale $= 1 + \sqrt{3} \approx 2{,}7321$.

\subsection{Partition $\{T_1, T_2\}|\{T_3, T_4\}$}
Par symétrie horizontale, la solution est $S_1 \approx (0{,}5, 0{,}211)$
et $S_2 \approx (0{,}5, 0{,}789)$.
Longueur totale $= 1 + \sqrt{3} \approx 2{,}7321$.

\subsection{Résultat}
\[
    \ell_{\min} \;=\; 1 + \sqrt{3} \;\approx\; 2{,}73205,
\]
obtenu avec 2 points de Steiner. C'est la valeur exacte connue pour le carré unité.

% ==================================================================
\section{Complexité et convergence}
% ==================================================================

\begin{center}
\begin{tabular}{lcc}
\toprule
\textbf{Phase} & \textbf{Complexité} & \textbf{Convergence} \\
\midrule
MST (Prim) & $O(n^2) = O(16)$ & Exacte \\
Fermat par Weiszfeld & $O(\text{max\_iter})$ & Linéaire garantie \\
2-Steiner Weiszfeld alterné & $O(4 \times \text{max\_iter})$ & Linéaire (convexe) \\
Sélection du minimum & $O(16)$ & — \\
\midrule
\textbf{Total} & $O(\text{max\_iter})$ & Globalement optimal \\
\bottomrule
\end{tabular}
\end{center}

\medskip

La limite de \textbf{50\,000 itérations} est largement suffisante en pratique :
la convergence est typiquement atteinte en quelques centaines d'itérations pour
des coordonnées d'écran (ordre de grandeur : pixels).

\end{document}
